\chapter{Réalisation du frontend}
\label{ch:frontend}

Toutes les captures de ce chapitre proviennent de l'application en fonctionnement,
connectée à la base de données réelle. Le jeu de données affiché est un jeu de
démonstration constitué pendant le développement ; aucune image n'est une maquette.

\section{Architecture React}

L'interface est une application monopage organisée en quatre niveaux. Les \emph{pages}
--- 21 au total --- correspondent aux routes et orchestrent le chargement des données. Les
\emph{composants} --- 27 --- encapsulent les fragments réutilisables, principalement les
formulaires modaux. Le \emph{système de composants}, rassemblé dans
\code{components/ui/}, fournit sept modules de primitives : boutons, badges, tables,
retours d'état, mise en page, graphiques. Les \emph{services} --- 19 modules --- isolent
tous les appels HTTP ; aucune page n'appelle directement Axios.

Cette séparation a une conséquence pratique vérifiée lors de l'audit : lorsqu'un contrat
d'API a changé, une seule fonction de service a dû être corrigée, et l'ensemble des pages
consommatrices en a bénéficié.

Le routage distingue trois espaces : les routes publiques, les routes d'authentification,
et les routes d'administration regroupées sous \code{/admin}. Ces dernières sont protégées
par un composant \code{RequireRole} qui redirige vers le tableau de bord tout utilisateur
dont le rôle ne figure pas dans la liste autorisée.

\section{Authentification}

L'écran de connexion (figure~\ref{fig:login}) est adossé à l'endpoint \code{/auth/login}.
Il gère l'état de chargement, l'erreur d'identifiants et le message de session expirée ou
de compte gelé, déposé par l'intercepteur Axios lorsqu'une réponse~401 survient en cours de
navigation.

\begin{figure}[H]
\centering
\includegraphics[width=\textwidth]{\figpath/screens/03-login.png}
\caption{Interface d'authentification}
\label{fig:login}
\end{figure}

Le jeton reçu est conservé côté navigateur et réinjecté par l'intercepteur de requête sur
chaque appel. L'écran donne également accès au parcours de réinitialisation du mot de
passe.

\section{Tableau de bord}

Le tableau de bord (figure~\ref{fig:dashboard}) est la page d'entrée de l'espace interne.
Son contenu est adapté au rôle : les indicateurs de demandes d'adhésion et les blocs de
pilotage ne sont chargés que pour les profils de direction, ce qui évite d'émettre des
requêtes vouées à un refus.

\begin{figure}[H]
\centering
\includegraphics[width=\textwidth]{\figpath/screens/04-dashboard.png}
\caption{Tableau de bord principal du système LIAS}
\label{fig:dashboard}
\end{figure}

On y distingue la barre de navigation latérale, groupée par domaine --- laboratoire,
activité scientifique, ressources, pilotage --- et la barre supérieure portant la recherche
globale et le compteur de notifications non lues.

\section{Gestion des membres}

La liste des membres (figure~\ref{fig:membres}) propose une recherche temporisée, un
filtrage par statut et des indicateurs de répartition. Le statut est rendu par un badge de
couleur distincte, conformément à l'exigence de lisibilité de la séparation entre statut et
rôle.

\begin{figure}[H]
\centering
\includegraphics[width=\textwidth]{\figpath/screens/05-membres-liste.png}
\caption{Liste des membres du laboratoire avec filtrage par statut}
\label{fig:membres}
\end{figure}

La sélection d'une ligne ouvre la fiche du membre, organisée en quatre onglets. L'onglet
\emph{Gouvernance} (figure~\ref{fig:membreGouv}) est celui qui matérialise la règle
centrale du projet.

\begin{figure}[H]
\centering
\includegraphics[width=\textwidth]{\figpath/screens/07-membre-gouvernance.png}
\caption{Fiche membre --- onglet Gouvernance, matérialisant la séparation \srt{} / \rol{} / \man{}}
\label{fig:membreGouv}
\end{figure}

Cet écran présente côte à côte le statut d'appartenance et le mandat de gouvernance en
cours, accompagnés d'un encadré explicitant que le rôle effectif est \emph{dérivé} des deux
et n'est jamais saisi. Sur la capture, le membre est doctorante, ne détient aucun mandat
actif, et l'historique conserve néanmoins un mandat de chef d'équipe échu : la séparation
des trois notions est directement lisible à l'écran, ce qui était l'un des objectifs
d'interface du cahier des charges.

L'onglet \emph{Historique} (figure~\ref{fig:membreHist}) présente en vis-à-vis les deux
chronologies exigées par la section~4 du cahier des charges : l'historique des statuts et
les périodes d'affiliation.

\begin{figure}[H]
\centering
\includegraphics[width=\textwidth]{\figpath/screens/08-membre-historique.png}
\caption{Fiche membre --- historique des statuts et périodes d'affiliation}
\label{fig:membreHist}
\end{figure}

La mention portée sous la seconde colonne --- un membre peut quitter puis revenir, chaque
période étant conservée séparément --- rappelle la règle §4.3 dont cet écran est la
traduction visuelle.

\section{Gouvernance}

Le module de gouvernance est organisé en trois onglets. L'onglet \emph{Structure}
(figure~\ref{fig:gouvStruct}) affiche l'organigramme de direction, entièrement déduit des
mandats actifs : un poste sans titulaire est explicitement présenté comme vacant plutôt que
masqué.

\begin{figure}[H]
\centering
\includegraphics[width=\textwidth]{\figpath/screens/09-gouvernance-structure.png}
\caption{Gouvernance --- organigramme de direction déduit des mandats actifs}
\label{fig:gouvStruct}
\end{figure}

L'onglet \emph{Mandats} (figure~\ref{fig:gouvMandats}) permet l'attribution et la clôture
des responsabilités. Le formulaire avertit que l'attribution d'un mandat de directeur ou de
vice-directeur clôt automatiquement celui du titulaire précédent, et la note de bas de
tableau rappelle qu'une clôture ne supprime rien.

\begin{figure}[H]
\centering
\includegraphics[width=\textwidth]{\figpath/screens/11-gouvernance-mandats.png}
\caption{Gouvernance --- attribution et clôture des mandats}
\label{fig:gouvMandats}
\end{figure}

Ces commandes ne sont affichées qu'à l'administrateur système et au directeur. Un
vice-directeur accédant à la même page obtient la consultation et un encart explicatif ---
correction issue de l'audit, décrite au chapitre~\ref{ch:tests}.

\section{Équipes}

L'onglet \emph{Équipes} (figure~\ref{fig:gouvEquipes}) présente les équipes de recherche
avec leur responsable, leur effectif et leur date de création, ainsi que le formulaire de
transfert d'un membre d'une équipe à une autre.

\begin{figure}[H]
\centering
\includegraphics[width=\textwidth]{\figpath/screens/10-gouvernance-equipes.png}
\caption{Gouvernance --- équipes de recherche et transfert de membre}
\label{fig:gouvEquipes}
\end{figure}

Le responsable affiché n'est pas un champ saisi : il provient du mandat de chef d'équipe
actif rattaché à l'équipe. Une équipe sans mandat actif porte la mention « Aucun chef
d'équipe ».

\section{Demandes d'adhésion}

L'écran des candidatures (figure~\ref{fig:demandes}) liste les dossiers déposés depuis le
portail public, avec leur état et le détail de chaque dossier.

\begin{figure}[H]
\centering
\includegraphics[width=\textwidth]{\figpath/screens/12-demandes-adhesion.png}
\caption{Demandes d'adhésion déposées depuis le portail public}
\label{fig:demandes}
\end{figure}

Conformément à la section~6 du cahier des charges, les boutons de décision ne sont rendus
que pour le directeur en mandat ; l'administrateur système et le vice-directeur accèdent à
la consultation, accompagnée d'un message explicite. L'acceptation crée le compte du
membre, ouvre sa première période d'affiliation et lui transmet ses identifiants.

\section{Équipements}

Le module matériel (figure~\ref{fig:equip}) est organisé en deux onglets : l'inventaire et
les demandes. Les indicateurs de tête distinguent les références en stock, les unités
disponibles, les unités distribuées et les références épuisées.

\begin{figure}[H]
\centering
\includegraphics[width=\textwidth]{\figpath/screens/13-equipements.png}
\caption{Gestion du matériel --- inventaire et indicateurs de stock}
\label{fig:equip}
\end{figure}

Le registre de distribution (figure~\ref{fig:equipDist}) répond à l'exigence la plus
spécifique de la section~14 : il présente non seulement les bénéficiaires d'un article,
mais également les membres actifs qui ne l'ont pas reçu, information calculée côté serveur.

\begin{figure}[H]
\centering
\includegraphics[width=\textwidth]{\figpath/screens/27-equipement-distribution.png}
\caption{Registre de distribution --- bénéficiaires et non-bénéficiaires d'un article}
\label{fig:equipDist}
\end{figure}

\section{Documents}

L'archivage documentaire (figure~\ref{fig:docs}) propose la recherche, le filtrage par
catégorie, le dépôt et le téléchargement, avec des indicateurs de volume stocké et de
documents versionnés.

\begin{figure}[H]
\centering
\includegraphics[width=\textwidth]{\figpath/screens/14-documents.png}
\caption{Archivage documentaire du laboratoire}
\label{fig:docs}
\end{figure}

L'historique des versions (figure~\ref{fig:docVersions}) matérialise l'exigence
d'historisation : chaque dépôt successif est conservé avec son numéro, son auteur et sa
date, et reste téléchargeable individuellement.

\begin{figure}[H]
\centering
\includegraphics[width=\textwidth]{\figpath/screens/26-document-versions.png}
\caption{Historique des versions d'un document archivé}
\label{fig:docVersions}
\end{figure}

Les téléchargements ne sont pas de simples liens : les fichiers étant hors de l'espace
public, ils sont récupérés par requête authentifiée puis remis au navigateur, ce qui
garantit que le contrôle d'accès s'applique aussi aux pièces jointes.

\section{Événements et calendrier}

L'écran des événements (figure~\ref{fig:events}) couvre les conférences, séminaires et
workshops, avec la gestion des organisateurs et des inscriptions.

\begin{figure}[H]
\centering
\includegraphics[width=\textwidth]{\figpath/screens/15-evenements.png}
\caption{Événements scientifiques du laboratoire}
\label{fig:events}
\end{figure}

Le calendrier (figure~\ref{fig:calendrier}) agrège en une vue unique les événements
scientifiques et les réunions internes, avec navigation mensuelle et panneau de détail du
jour sélectionné.

\begin{figure}[H]
\centering
\includegraphics[width=\textwidth]{\figpath/screens/16-calendrier.png}
\caption{Calendrier unifié des événements et des réunions}
\label{fig:calendrier}
\end{figure}

Cette agrégation est réalisée côté serveur par \code{CalendarService}, qui fusionne deux
sources hétérogènes en une liste ordonnée d'éléments typés.

\section{Publications et procès-verbaux}

Le répertoire des publications (figure~\ref{fig:publis}) permet la recherche et le
classement par année, équipe et auteur, conformément à la section~12. Son ajout est ouvert
aux doctorants comme aux membres.

\begin{figure}[H]
\centering
\includegraphics[width=\textwidth]{\figpath/screens/17-publications.png}
\caption{Répertoire des publications scientifiques}
\label{fig:publis}
\end{figure}

Le module des réunions (figure~\ref{fig:reunions}) couvre la planification, l'ordre du jour
et le dépôt des procès-verbaux exigés par la section~17.

\begin{figure}[H]
\centering
\includegraphics[width=\textwidth]{\figpath/screens/18-reunions-pv.png}
\caption{Réunions internes et archivage des procès-verbaux}
\label{fig:reunions}
\end{figure}

Le module des conventions (figure~\ref{fig:conventions}) gère les partenaires académiques
et industriels ainsi que les conventions associées et leurs documents.

\begin{figure}[H]
\centering
\includegraphics[width=\textwidth]{\figpath/screens/25-conventions.png}
\caption{Conventions et partenariats du laboratoire}
\label{fig:conventions}
\end{figure}

\section{Notifications}

L'espace de communication interne (figure~\ref{fig:notifs}) réunit le centre de
notifications personnel et les discussions, déclinées selon les trois portées prévues par
la section~13 : laboratoire, équipe et événement.

\begin{figure}[H]
\centering
\includegraphics[width=\textwidth]{\figpath/screens/19-notifications.png}
\caption{Notifications personnelles et espaces de discussion}
\label{fig:notifs}
\end{figure}

Un compte technique dépourvu de profil de membre --- l'administrateur système --- ne
possède pas de boîte de réception ; l'interface le détecte et affiche un état vide explicite
plutôt qu'une erreur.

\section{Audit et rapport annuel}

Le journal d'audit (figure~\ref{fig:audit}) restitue la trace intégrale des actions
sensibles : acteur, action, entité concernée, adresse IP et horodatage.

\begin{figure}[H]
\centering
\includegraphics[width=\textwidth]{\figpath/screens/20-audit.png}
\caption{Journal d'audit --- traçabilité des actions du système}
\label{fig:audit}
\end{figure}

L'écran des rapports d'activité (figure~\ref{fig:rapports}) permet de générer la synthèse
d'une année et d'en télécharger la version PDF.

\begin{figure}[H]
\centering
\includegraphics[width=\textwidth]{\figpath/screens/21-rapports.png}
\caption{Rapports d'activité annuels et export PDF}
\label{fig:rapports}
\end{figure}

L'administration système (figure~\ref{fig:admin}) regroupe la gestion des comptes, des
rôles techniques et de l'identité institutionnelle. L'avertissement porté en bas d'écran
rappelle qu'une modification manuelle de rôle constitue une dérogation administrative,
historisée au même titre qu'un changement issu d'un mandat.

\begin{figure}[H]
\centering
\includegraphics[width=\textwidth]{\figpath/screens/22-administration.png}
\caption{Administration système --- comptes, rôles et identité du laboratoire}
\label{fig:admin}
\end{figure}

\section{Portail public}

Le portail public (figure~\ref{fig:portail}) constitue la vitrine institutionnelle. Il
présente le laboratoire, ses indicateurs, ses publications récentes, ses équipes et ses
événements, sans exiger d'authentification.

\begin{figure}[H]
\centering
\includegraphics[width=\textwidth]{\figpath/screens/01-portail-public.png}
\caption{Portail public du laboratoire LIAS}
\label{fig:portail}
\end{figure}

Ses données proviennent des mêmes tables que l'espace interne, via des endpoints
explicitement ouverts. Le formulaire de candidature (figure~\ref{fig:apply}) permet le
dépôt d'un dossier avec lettre de motivation et curriculum vitæ, alimentant directement le
module de demandes d'adhésion.

\begin{figure}[H]
\centering
\includegraphics[width=\textwidth]{\figpath/screens/02-portail-apply.png}
\caption{Formulaire public de candidature au laboratoire}
\label{fig:apply}
\end{figure}

\section{Design responsive et navigation}

La mise en page repose sur les grilles et utilitaires de Tailwind. Sous 1024~pixels, la
barre latérale bascule en tiroir superposé avec voile de fond ; les grilles d'indicateurs
passent de quatre à deux puis à une colonne ; les tableaux larges défilent horizontalement
dans leur propre conteneur, de sorte que la page elle-même ne défile jamais latéralement.

La navigation latérale est générée à partir d'une description déclarative associant à
chaque entrée la liste des rôles autorisés, identique à celle qu'applique le garde de
route. Un module inaccessible n'est donc pas affiché puis refusé : il est absent de la
navigation. Cette propriété est illustrée au chapitre suivant
(figure~\ref{fig:sidebarPhd}).
